% =============================================================================
% 4.2 Admission. Ben, 2026-10-18: admission is the EXECUTION of the capacity
% model's decision (s3); it classifies requests to fill the optimal rider count.
% Which requests: LARGEST FIRST (review note 10) -- context cancels in the
% workload view, not in admission. Where the code (vllm/v1/core/sched/
% transfer_aware_admission.py on perblock-d2h-opt) does not yet do what the
% algorithm says, the gap is listed in Appendix C (sections/C-wip.tex),
% never in the body (Ben, 2026-10-19).
% =============================================================================

\subsection{Admission}\label{sec:admission}

The admission controller carries out the capacity model's decision
(\cref{eq:decision}) every step. It keeps $T_0$ current, re-derives $\rmin$
from it, and assigns each request a tier, adding riders while the step-time
law of \cref{eq:step} predicts that each one raises throughput
(\cref{alg:admit}).

\paragraph{Which requests ride: the largest.} When more requests are waiting
than HBM can hold, the controller must decide which of them take HBM and which
take the link. Each request spends its bytes from exactly one budget: HBM
capacity as a resident, or link bytes per step as a rider. Whichever budget it
spends, it delivers one token per step. The tiers differ in what else they charge. A resident costs
nothing beyond its bytes, while every rider also costs $c$ of step time,
however long its context. A large rider spends that $c$ on many hidden bytes,
and a small one on few. Taking the largest candidates first therefore
moves the most bytes for the least added step time. It also leaves the smallest
requests for HBM, which is the larger budget and so holds the most of them. A
cohort of large riders is also a small one, so the controller still checks
$\rmin$. At our operating point $\rmin = 1$, so this check does not bind.

\paragraph{Keeping $T_0$ honest.} The prediction starts from $T_0$, and a step
that already carries riders is slower. Measuring $T_0$ on such steps would
count the riders' own cost as room for more riders, which would admit more
and slow the step further. $T_0$
is therefore an average over steps that carry no rider (line~\ref{alg:t0}).

\paragraph{Every migration is a hypothesis.} A move the model predicts to pay can
still be wrong for latency. On A100, a demotion
that passed the byte test raised inter-token latency from about 6 to 24\,ms. The
controller therefore moves running requests between tiers provisionally. It
records the best recent step time, demotes the largest resident whose move the
model predicts pays, waits for the step time to settle, and undoes the move if latency rose.
Each undo doubles a back-off, so the same move is not retried at once.
% src: transfer_aware_admission.py: hide_budget_ms (rider-free EMA),
%      plan_migrations (probe, min-of-window baseline, exponential backoff);
%      the 6 -> 24 ms regression is recorded in its comments.

\begin{algorithm}[t]
\caption{Admission: executing the capacity model's decision, once per
scheduling step.}
\label{alg:admit}
\Input{profile $(F, c, \delta, \Xlink, \btok)$; KV memory $\Mkv$;
  waiting queue $\mathcal{W}$; running set $\mathcal{R}$}
\Output{a tier for every request in $\mathcal{R}\cup\mathcal{W}$}
\BlankLine
\Proc{\textsc{Admit}$(\mathcal{W}, \mathcal{R})$}{
  $T_0 \leftarrow$ mean step time over rider-free steps\;\label{alg:t0}
  $\rmin \leftarrow$ capacity model at $T_0$\cmt{Eq. \ref{eq:decision}}\label{alg:model}
  $\mathit{bytes} \leftarrow \btok \sum_{i \in \mathrm{riders}(\mathcal{R})} S_i$;\quad
  $n \leftarrow |\mathrm{riders}(\mathcal{R})|$;\quad
  $\mathcal{N} \leftarrow \emptyset$\;
  \ForEach{$q \in \mathcal{W}$ \textnormal{by context, smallest first}}{\label{alg:res}
    \lIf{$q$ \textnormal{fits in free HBM}}{$\mathit{tier}(q) \leftarrow$ \textsf{resident}}
  }
  \ForEach{\textnormal{remaining} $q \in \mathcal{W}$ \textnormal{by context, largest first}}{\label{alg:order}
    \eIf{$n + 1 \ge \rmin$ \textnormal{and}
         $\Theta(\mathit{bytes} + S_q\btok,\, n+1) \le \Theta(\mathit{bytes}, n)$}{\label{alg:knee}
      defer $q$\cmt{past the knee, no longer pays}
    }{
      $\mathcal{N} \leftarrow \mathcal{N}\cup\{q\}$;\quad
      $\mathit{bytes} \leftarrow \mathit{bytes} + S_q\btok$;\quad
      $n \leftarrow n+1$\;
    }
  }
  \eIf{$|\mathrm{riders}(\mathcal{R})| + |\mathcal{N}| < \rmin$}{\label{alg:commit}
    defer all of $\mathcal{N}$\cmt{cannot repay F}
  }{
    $\mathit{tier}(q) \leftarrow$ \textsf{rider} for all $q \in \mathcal{N}$\;
  }
}
\BlankLine
\Proc{\textsc{Migrate}$(\mathcal{R})$}{
  \If{\textnormal{a probe on $d$ has settled}}{
    \If{$\mathrm{ITL} > (1+\epsilon)\cdot\mathit{base}$}{
      promote $d$;\quad $\mathit{backoff} \leftarrow 2\,\mathit{backoff}$\;
    }
    \Return\;
  }
  \If{\textnormal{HBM is full and $\mathcal{W}$ is not empty}}{\label{alg:migr}
    $d \leftarrow$ largest resident whose demotion $\Theta$ predicts pays\;
    $\mathit{base} \leftarrow \min(\textnormal{recent ITL})$\cmt{not the mean}
    demote $d$; open a probe on $d$\;
  }
}
\end{algorithm}

\Cref{alg:admit} states the controller. $\Theta(\mathit{bytes}, n)$ is the
predicted throughput, $(B + n)/T$ with $T$ from \cref{eq:step} for $n$ riders
carrying $\mathit{bytes}$ ($T = T_0$ when $n = 0$). \textsc{Admit} places new
requests in two passes: residents smallest first while HBM has room, then
riders largest first while each still raises $\Theta$. The riders are committed only if there are at least $\rmin$ of
them.
\textsc{Migrate} demotes the largest resident when HBM is full and requests are
waiting, which makes room for smaller ones, and it treats every move as
provisional.
